\begin{theorem}[square-zero positive ideal]
\label{thm:bcpr-square-zero-algebra}
\Verified\ Let $N\subset H^{*,*}_{BC}(X)$ be the span of the nine positive-degree
generators in \cref{eq:bcpr-generators}.  Then
\begin{equation}
 H^{*,*}_{BC}(X)=\mathbb C\mathbf1\oplus N,
 \qquad N^2=0.
\label{eq:bcpr-algebra}
\end{equation}
Equivalently, every product of two positive-degree Bott--Chern classes is
zero.
\end{theorem}

\begin{proof}
We first calculate the products landing below top degree.  Since
$\bar\partial G=\bar\partial Q=0$ and, after conjugation,
$\partial\bar G=\partial\bar Q=0$, the graded Leibniz rule gives the exact
identities
\begin{equation}
 \partial\bar\partial(G\wedge\bar G)
 =-a\wedge a^\dagger,
\label{eq:bcpr-product-ab}
\end{equation}
\begin{equation}
 \partial\bar\partial(G\wedge\bar Q)
 =-a\wedge c^\dagger,
 \qquad
 \partial\bar\partial(\bar G\wedge Q)
 =a^\dagger\wedge c.
\label{eq:bcpr-product-acbar-bc}
\end{equation}
Consequently
\begin{equation}
 \mathsf A\mathsf B=0,\qquad
 \mathsf A\overline{\mathsf C}=0,\qquad
 \mathsf B\mathsf C=0.
\label{eq:bcpr-first-zero-products}
\end{equation}

Choose the pullback Fubini--Study metric on
$f^*\mathcal O_B(1)$.  Its Chern form is the pullback of a $(1,1)$ form on
the curve $B$, so its square vanishes pointwise.  Squaring
\cref{eq:bcpr-c1-pullback} and using
\cref{eq:bcpr-first-zero-products} gives
\begin{equation}
 \mathsf A^2+\mathsf B^2=0.
\label{eq:bcpr-sum-of-squares}
\end{equation}
\Cref{eq:bcpr-sum-of-squares} alone does not separate the two
squares.  Apply $\pi_{\bar\partial}^{2,2}$.  The class
$\mathsf B^2$ has zero Dolbeault image because
\begin{equation}
 (a^\dagger)^2
 =\bar\partial(\bar G\wedge a^\dagger).
\label{eq:bcpr-b-square-dolbeault-zero}
\end{equation}
It follows from \cref{eq:bcpr-sum-of-squares} that
$[a^2]_{\bar\partial}=0$.  On the other hand,
\begin{equation}
 [a^2]_{\bar\partial}
 =d_1[G\wedge a]_{\bar\partial},
 \qquad
 \partial(G\wedge a)=a^2.
\label{eq:bcpr-a-square-d1}
\end{equation}
The first arrow in \cref{eq:bcpr-hodge-row-two} is injective, so
\cref{eq:bcpr-a-square-d1} implies
$[G\wedge a]_{\bar\partial}=0$.  There is therefore a form
$\Phi\in\mathcal A^{1,1}(X)$ such that
\begin{equation}
 G\wedge a=\bar\partial\Phi.
\label{eq:bcpr-phi-primitive}
\end{equation}
Applying $\partial$ gives
\begin{equation}
 a^2=\partial\bar\partial\Phi,
\label{eq:bcpr-a-square-bc-zero}
\end{equation}
and hence $\mathsf A^2=0$.  Conjugation gives
$\mathsf B^2=0$.

It remains among the products below top bidegree to prove
$\mathsf A\mathsf C=0$; its conjugate will give
$\mathsf B\overline{\mathsf C}=0$.  First, the natural map
\begin{equation}
 \pi_{\bar\partial}^{3,2}\colon H^{3,2}_{BC}(X)
 \longrightarrow H^{3,2}_{\bar\partial}(X)
\label{eq:bcpr-bc-to-dolbeault-32}
\end{equation}
is an isomorphism.  Indeed, both spaces have dimension one by
\cref{eq:bcpr-nonzero-dimensions,cor:exact-hodge-diamond}, and
\cref{eq:bcpr-e-nonzero} says that $\mathsf E$ has nonzero image.
Pair $\pi_{\bar\partial}^{3,2}(\mathsf A\mathsf C)$ with the generator
$[G]_{\bar\partial}\in H^{0,1}_{\bar\partial}(X)$.  By
\cref{eq:bcpr-phi-primitive}, $\bar\partial c=0$, and Stokes' theorem,
\begin{equation}
 \int_Xa\wedge c\wedge G
 =-\int_XG\wedge a\wedge c
 =-\int_X\bar\partial\Phi\wedge c
 =-\int_X\bar\partial(\Phi\wedge c)=0.
\label{eq:bcpr-serre-pairing-zero}
\end{equation}
Dolbeault--Serre duality makes
\begin{equation}
 H^{3,2}_{\bar\partial}(X)\times
 H^{0,1}_{\bar\partial}(X)\longrightarrow\mathbb C,
 \qquad
 ([\alpha],[\rho])\longmapsto\int_X\alpha\wedge\rho,
\label{eq:bcpr-serre-pairing}
\end{equation}
a perfect pairing.  Since
$H^{0,1}_{\bar\partial}(X)=\mathbb C[G]_{\bar\partial}$,
\cref{eq:bcpr-serre-pairing-zero} implies
$\pi_{\bar\partial}^{3,2}(\mathsf A\mathsf C)=0$; injectivity in
\cref{eq:bcpr-bc-to-dolbeault-32} gives
\begin{equation}
 \mathsf A\mathsf C=0,
 \qquad
 \mathsf B\overline{\mathsf C}=0.
\label{eq:bcpr-second-zero-products}
\end{equation}
This step uses Dolbeault--Serre duality together with the independently
proved injectivity of \cref{eq:bcpr-bc-to-dolbeault-32}.  The
Bott--Chern--Aeppli pairing of
\cite[\S2.c]{Schweitzer2007} by itself would supply a perfect pairing but
would not annihilate this product.

The only products not yet considered which can have bidegree at most
$(3,3)$ are
\begin{equation}
 \mathsf A\mathsf D,\quad
 \mathsf A\overline{\mathsf D},\quad
 \mathsf B\mathsf D,\quad
 \mathsf B\overline{\mathsf D},\quad
 \mathsf C\overline{\mathsf C},
\label{eq:bcpr-top-products}
\end{equation}
all of bidegree $(3,3)$.  Every non-top generator in
\cref{eq:bcpr-generators} is de Rham exact:
\begin{equation}
\begin{gathered}
 a=dG,\qquad a^\dagger=d\bar G,\qquad
 c=dQ,\qquad c^\dagger=d\bar Q,\\
 d_0=d(G\wedge Q),\qquad
 d_0^\dagger=d(\bar G\wedge\bar Q),\qquad
 e=d(Q^2),\qquad e^\dagger=d(\bar Q^2).
\end{gathered}
\label{eq:bcpr-de-rham-primitives}
\end{equation}
Each factor paired with an exact factor is $d$-closed, so every product in
\cref{eq:bcpr-top-products} is $d$-exact and has integral zero.  Integration
induces an isomorphism
\begin{equation}
 H^{3,3}_{BC}(X)\xrightarrow{\ \int_X\ }\mathbb C.
\label{eq:bcpr-top-integration}
\end{equation}
It is well defined by Stokes' theorem, the source has dimension one by
\cref{eq:bcpr-nonzero-dimensions}, and it is nonzero on $\mathsf T$.
Therefore all five products in \cref{eq:bcpr-top-products} vanish in
Bott--Chern cohomology.

All other products of positive-degree generators have $p>3$ or $q>3$ and
vanish already as differential forms.
\Cref{eq:bcpr-first-zero-products,eq:bcpr-a-square-bc-zero,eq:bcpr-second-zero-products,eq:bcpr-top-products,eq:bcpr-top-integration}
exhaust every possible positive-degree product
and prove \cref{eq:bcpr-algebra}.
\end{proof}
